\DocumentMetadata{
  pdfversion=2.0,
  pdfstandard=ua-2,
  lang=en-US,
  tagging=on,
  tagging-setup={math/setup=mathml-SE,tabsorder=structure}
}
\documentclass[11pt,letterpaper]{article}
\usepackage[margin=0.68in,includefoot]{geometry}
\usepackage{newcomputermodern}
\usepackage{amsmath,amscd,mathtools}
\usepackage{tikz,adjustbox}
\providecommand{\square}{\mdlgwhtsquare}
\usepackage[hidelinks]{hyperref}
\hypersetup{
  pdftitle={Analysis PhD Exam, May 2002},
  pdfauthor={University of Florida Department of Mathematics},
  pdfsubject={Graduate mathematics examination},
  pdfdisplaydoctitle=true,
  bookmarksopen=true
}
\setcounter{secnumdepth}{0}
\setlength{\parindent}{0pt}
\setlength{\parskip}{0.28em}
\setlength{\emergencystretch}{3em}
\setlength{\leftmargini}{2.15em}
\setlength{\leftmarginii}{2.65em}
\setlength{\leftmarginiii}{3em}
\pagestyle{empty}
\raggedbottom
\allowdisplaybreaks
\NewTaggingSocketPlug{block/list/label}{examproblem}{%
  \tagstructbegin{tag=Lbl}\tagstructend%
  \tagstructbegin{tag=\LBody}%
  \tagstructbegin{tag=\ProblemHeadingTag,title-o={\ProblemHeadingText}}%
  \tagmcbegin{tag=\ProblemHeadingTag,actualtext={\ProblemHeadingText}}%
  #2%
  \tagmcend%
  \tagstructend%
}
\newenvironment{examproblems}[2]{%
  \def\ProblemHeadingTag{#2}%
  \AssignTaggingSocketPlug{block/list/label}{examproblem}%
  \begin{enumerate}%
  \setcounter{enumi}{#1}%
  \renewcommand{\labelenumi}{\textbf{\theenumi.}}%
  \setlength{\topsep}{0.35em}%
  \setlength{\partopsep}{0pt}%
  \setlength{\itemsep}{0.48em}%
  \setlength{\parsep}{0.18em}%
}{\end{enumerate}}
\newenvironment{examsubparts}{%
  \AssignTaggingSocketPlug{block/list/label}{default}%
  \begin{enumerate}%
  \setlength{\topsep}{0.18em}%
  \setlength{\partopsep}{0pt}%
  \setlength{\itemsep}{0.16em}%
  \setlength{\parsep}{0.1em}%
}{\end{enumerate}}
\newenvironment{examinstructions}{%
  \par\begingroup\itshape
}{%
  \par\endgroup\vspace{0.18em}
}
\makeatletter
\renewcommand\subsection{\@startsection{subsection}{2}{0pt}%
  {-1.25ex plus -.25ex minus -.1ex}%
  {0.55ex plus .1ex}%
  {\normalfont\large\bfseries}}
\renewcommand{\@maketitle}{%
  \newpage\null\vskip -1em%
  {\footnotesize\bfseries
    \href{https://www.ufl.edu/}{University of Florida}, \href{https://math.ufl.edu/}{Department of Mathematics}\par}%
  \vskip -0.5em%
  \tagstructbegin{tag=H1,title={Analysis PhD Exam, May 2002}}%
  \tagpdfparaOff%
  \tagmcbegin{tag=H1}%
    {\LARGE\bfseries\href{https://gma.math.ufl.edu/exams/analysis-phd/}{Analysis PhD Exam}, May 2002\par}%
  \tagmcend%
  \tagpdfparaOn%
  \tagstructend%
  \vskip 0.25em%
  \tagmcbegin{artifact}%
    \hrule height 1pt%
  \tagmcend%
  \vskip 1em%
}
\makeatother
\title{Analysis PhD Exam, May 2002}
\author{}
\date{}
\begin{document}
\maketitle
\thispagestyle{empty}
\subsection{I}
\begin{examinstructions}
State the following theorems:
\end{examinstructions}
\begin{examproblems}{0}{H3}
\def\ProblemHeadingText{Problem 1.}
\item
The Egorov theorem.
\def\ProblemHeadingText{Problem 2.}
\item
The Radon–Nikodym theorem.
\def\ProblemHeadingText{Problem 3.}
\item
Density of step functions in \(L^p\) and \(L^\infty\).
\def\ProblemHeadingText{Problem 4.}
\item
The Fubini theorem for \(\mu\times\nu\)-integrable functions.
\def\ProblemHeadingText{Problem 5.}
\item
The Riesz representation theorem of the dual of \(C(K)\).
\end{examproblems}
\subsection{II}
\begin{examinstructions}
State and prove one of the following theorems:
\end{examinstructions}
\begin{examproblems}{0}{H3}
\def\ProblemHeadingText{Problem 1.}
\item
The integral representation of the dual of \(L^1(\mu)\) (the finite measure case).
\def\ProblemHeadingText{Problem 2.}
\item
The theorem of integration with respect to \(g\mu\).
\end{examproblems}
\subsection{III}
\begin{examinstructions}
Solve the following problems:
\end{examinstructions}
\begin{examproblems}{0}{H3}
\def\ProblemHeadingText{Problem 1.}
\item
Let \((X,\Sigma,\mu)\) be a measure space with \(\mu(X)<\infty\), \(F\) a Banach space, and \(\mathcal F\subseteq\Sigma\) a~\(\sigma\)-algebra. Prove the existence of the conditional expectation \(E(f\mid\mathcal F)\) for \(f\in L_F^1(\mu)\).
\def\ProblemHeadingText{Problem 2.}
\item
Let \((X,\mathcal S)\), \((Y,\mathcal T)\) be measurable spaces and \(\mu:\mathcal S\to\mathbb R\), \(\nu:\mathcal T\to\mathbb R\) σ-additive measures. Prove the existence of the product measure \(\mu\times\nu\) and the equality \(\lVert\mu\times\nu\rVert=\lVert\mu\rVert\times\lVert\nu\rVert\).
\def\ProblemHeadingText{Problem 3.}
\item
Let \((X,\mathcal S)\), \((Y,\mathcal T)\) be measurable spaces, \(p:X\to Y\) an \((\mathcal S,\mathcal T)\)-measurable mapping, \(\mu:\mathcal S\to\mathbb R_+\) a σ-additive measure, and \(\nu=p(\mu)\) the image measure of~\(\mu\) under the mapping~\(p\). Prove that if \(f\in L_F^1(\nu)\), then \(f\circ p\in L_F^1(\mu)\) and
\[
\int f\circ p\,d\mu=\int f\,dp(\mu).
\]
\def\ProblemHeadingText{Problem 4.}
\item
Let \((X,\Sigma,\mu)\) be a measure space and \(f\in L^1(\mu)\). Prove that for every \(\varepsilon>0\), there is a set \(X_\varepsilon\in\Sigma\) with \(\mu(X_\varepsilon)<\infty\) such that
\[
\int_{X\setminus X_\varepsilon}|f|\,d\mu<\varepsilon.
\]
\def\ProblemHeadingText{Problem 5.}
\item
Let~\(\lambda\) be the Lebesgue measure on \([0,1]\) and \(f\in L^1(\lambda)\). Prove that if
\[
\int_{[0,1]}x^n f\,d\lambda=0
\]
 for every~\(n\), then \(f=0\), \(\lambda\)-a.e.

\par
Hint: Use the fact that the continuous functions on \([0,1]\) are dense in \(L^1(\lambda)\) and reduce to the case \(\int_A f\,d\lambda=0\) for \(A\in\mathcal B([0,1])\).
\end{examproblems}
\begin{examinstructions}
Note. Write complete computations and explanations. Do not use arrows or other unnecessary signs. Use words for explanations. Write complete sentences. State the complete framework for each theorem.
\end{examinstructions}
\end{document}
